\section{Formal Methodology}
\label{sec:methodology}

Edge-deployed intrusion detection systems frequently require runtime interchangeability: a high-throughput primary engine (\eg ONNX Runtime) operating under steady state, and an in-process, zero-dependency fallback engine (\eg a pure NumPy vectorizer) engaged during container cold-starts, failovers, or constrained recovery phases. However, maintaining verdict equivalence across disparate execution engines is hindered by floating-point precision mismatches and threshold binning discretization. This section formalizes our two core methodological contributions: (i) boundary-directed parity testing to guarantee verdict consistency without requiring production traffic replay, and (ii) cache-blocked lockstep traversal with exact early exit to eliminate Python interpreter dispatch overhead.

\subsection{Decision Ensemble Formulation and the Precision Boundary}

Consider an ensemble $\mathcal{F} = \{T_1, \dots, T_M\}$ comprising $M$ decision trees trained on $d$-dimensional network flow features. Input vectors $\mathbf{x} \in \mathbb{R}^d$ are normalized via an affine transformation:
\begin{equation}
\mathbf{z} = (\mathbf{x} - \mathbf{o}) \odot \mathbf{s},
\end{equation}
where $\mathbf{o} \in \mathbb{R}^d$ denotes the feature centering offset, and $\mathbf{s} \in \mathbb{R}^d$ denotes the reciprocal feature scale vector.

Each internal node $v \in T_m$ is parameterized by a split feature index $g(v) \in \{1, \dots, d\}$ and a scalar threshold $\theta(v) \in \mathbb{R}$. During sample evaluation, routing branches left if $z_{g(v)} \le \theta(v)$ and right otherwise. 

\begin{definition}[$\epsilon$-Boundary Neighborhood]
Given a split node $v$ with threshold $\theta(v)$, the critical boundary neighborhood in raw feature space is defined by the interval:
\begin{equation}
\mathcal{B}(v) = \left[ x^* - \epsilon, x^* + \epsilon \right], \quad \text{where } x^* = \frac{\theta(v)}{s_{g(v)}} + o_{g(v)},
\end{equation}
and $\epsilon = 1\,\text{ulp}(x^*)$ denotes one unit in the last place of IEEE-754 single-precision representation.
\end{definition}

In standard replayed network traffic, the probability that an observed flow record falls into $\mathcal{B}(v)$ is bounded by the Lebesgue measure $\mu(\mathcal{B}(v)) \to 0$. Consequently, replayed datasets almost never stimulate borderline threshold discrepancies, masking divergence between float32 SIMD inference (\eg ONNX Runtime) and float64 CPU routines.

\begin{figure*}[t]
  \centering
  \resizebox{0.92\textwidth}{!}{%
  \begin{tikzpicture}[
      font=\small, >=Stealth, every node/.style={align=center},
      tn/.style ={draw=gray!60, thick, rounded corners=2pt, fill=gray!8, minimum width=1.55cm, minimum height=0.55cm, inner sep=2pt, font=\footnotesize},
      tgt/.style={tn, draw=accBlue, fill=accBlue!8, text=accBlue!75!black},
      leaf/.style={draw=gray!50, fill=white, circle, minimum size=0.32cm, inner sep=0pt},
      edge/.style={draw=gray!60, thick},
      path/.style={draw=accBlue, line width=1.4pt},
      lbl/.style={font=\footnotesize, inner sep=1.5pt, align=center},
      hdr/.style={font=\footnotesize\bfseries, align=left, anchor=west},
      probe/.style={circle, fill=accOrange, draw=white, line width=0.6pt, minimum size=0.26cm, inner sep=0pt},
      arr/.style ={-{Stealth[length=2.2mm]}, thick, draw=gray!60!black},
  ]
    % ---------------- (i) route a real row to the node ----------------
    \node[hdr] at (-0.4, 3.0) {(i) Route a real row to the node};
    \node[tn]  (r)  at (0.9, 2.3)  {$x_3 \le t_1$};
    \node[tn]  (a)  at (0.1, 1.3)  {$x_7 \le t_2$};
    \node[leaf](l1) at (1.7, 1.3)  {};
    \node[leaf](l2) at (-0.7, 0.3) {};
    \node[tgt] (n)  at (0.9, 0.3)  {$x_f \le t$};
    \node[leaf](l3) at (0.3, -0.6) {};
    \node[leaf](l4) at (1.5, -0.6) {};
    \draw[edge] (r) -- (l1);  \draw[edge] (a) -- (l2);
    \draw[edge] (n) -- (l3);  \draw[edge] (n) -- (l4);
    \draw[path] (r) -- (a);   \draw[path] (a) -- (n);
    \node[lbl, text=gray!60!black] at (-0.05, 1.95) {yes};
    \node[lbl, text=gray!60!black] at (0.85, 0.85)  {no};
    \node[lbl, anchor=west, text=accBlue!75!black] at (1.85, 0.3) {target node\\(one of 95{,}741)};
    % path constraints
    \node[lbl, anchor=north, draw=gray!40, rounded corners=2pt, fill=gray!4, inner sep=4pt, text width=4.2cm]
      (cons) at (0.9, -1.0)
      {intersect ancestor splits:\\
       $x_3 \in (-\infty, t_1]$, \quad $x_7 \in (t_2, \infty)$\\[1pt]
       write one value from each interval\\into a replay row $\Rightarrow$ row reaches $n$\\
       (99.9\% of nodes reached)};
    % ---------------- (ii) three float32 probes at the threshold ----------------
    \node[hdr] at (5.3, 3.0) {(ii) Overwrite $x_f$ with three float32 values};
    % scaled axis
    \draw[arr] (5.4, 1.4) -- (11.6, 1.4) node[lbl, anchor=west] {scaled $x_f$};
    \draw[thick, accBlue] (8.5, 1.1) -- (8.5, 1.75);
    \node[lbl, text=accBlue!75!black, anchor=south] at (8.5, 1.78) {threshold $t$};
    \node[probe] (p1) at (7.3, 1.4) {};
    \node[probe] (p2) at (8.5, 1.4) {};
    \node[probe] (p3) at (9.7, 1.4) {};
    \node[lbl, anchor=north] at (7.3, 1.2) {$t - 1$\,ulp};
    \node[lbl, anchor=north] at (8.5, 1.2) {$= t$};
    \node[lbl, anchor=north] at (9.7, 1.2) {$t + 1$\,ulp};
    \draw[gray!60, dashed] (8.5, 1.1) -- (8.5, 0.45);
    \node[lbl, anchor=south, text=gray!60!black] at (6.55, 1.55) {go left: $x_f \le t$};
    \node[lbl, anchor=south, text=gray!60!black] at (10.55, 1.55) {$x_f > t$: go right};
    % raw -> scaler -> scaled
    \node[lbl, anchor=south, text=gray!60!black] at (6.1, 2.1) {raw float32 $\xrightarrow{\;\text{scaler}\;}$ scaled};
    % float32 vs float64 comparison
    \node[lbl, anchor=north, draw=accGreen!80!black, rounded corners=2pt, fill=accGreen!8, inner sep=4pt, text width=3.1cm]
      at (6.6, 0.35)
      {\textbf{float32 scaler + compare}\\(ONNX Runtime, deployed NumPy)\\probe lands where placed:\\\textbf{0 flips} / 287{,}223};
    \node[lbl, anchor=north, draw=accOrange, rounded corners=2pt, fill=accOrange!8, inner sep=4pt, text width=3.1cm]
      at (10.5, 0.35)
      {\textbf{float64 scaler}\\scaled value differs by $\varepsilon$,\\the on-threshold probe can\\cross $t$: \textbf{328 flips}};
    \draw[arr, accOrange] (9.6, 0.35) to[out=105, in=-20] (8.72, 1.4);
  \end{tikzpicture}}
  \caption{Boundary-directed probe generation: (i)~path-satisfying routing to target internal nodes, and (ii)~threshold $\pm\,1$\,ulp float32 evaluation to detect precision-induced label flips.}
  \label{fig:probe}
\end{figure*}

\subsection{Algorithm 1: Boundary-Directed Probe Generation}

To expose rounding vulnerabilities deterministically without collecting millions of live flows, Algorithm~\ref{alg:boundary_probe} constructs synthetic boundary probes directly from the topology of the forest artifact, illustrated in Fig.~\ref{fig:probe}.

For every internal split node $v \in T_m$, the algorithm climbs backward along the unique path connecting $v$ to root $\text{root}(T_m)$. Each ancestor split $p$ restricts the feasible range $[L_f, U_f]$ of feature $f = g(p)$. By calculating the intersection across all ancestors, Algorithm~\ref{alg:boundary_probe} forms a feasible hyper-rectangle:
\begin{equation}
\mathcal{H}(v) = \prod_{f=1}^d [L_f, U_f],
\end{equation}
guaranteeing that any input vector drawn from $\mathcal{H}(v)$ will traverse unhindered from the root to node $v$. 

Once the feasible baseline $\mathbf{x}_{\text{base}}$ is formed, feature $g(v)$ is systematically perturbed across the three critical boundary conditions:
\begin{equation}
x' \in \{ x^*, \text{prev}_{32}(x^*), \text{next}_{32}(x^*) \}.
\end{equation}
This yields exactly $3 \cdot |\mathcal{V}_{\text{internal}}|$ synthetic probe rows. Evaluating this probe set across candidate runtimes provides an exact empirical upper bound on decision divergence.

\begin{algorithm}[t]
\caption{Boundary-Directed Probe Generation}
\label{alg:boundary_probe}
\footnotesize
\begin{algorithmic}[1]
\Require Forest $\mathcal{F} = \{T_1, \dots, T_M\}$, offset $\mathbf{o}$, scale $\mathbf{s}$
\Ensure Probe dataset $\mathcal{P} = \{(\mathbf{x}, v)\}$
\State $\mathcal{P} \gets \emptyset$
\ForAll{$T_m \in \mathcal{F}$}
    \State $\mathcal{V}_{\text{internal}} \gets \{v \in T_m \mid \neg \text{is\_leaf}(v)\}$
    \ForAll{$v \in \mathcal{V}_{\text{internal}}$}
        \LineComment{Step 1: Ascend to root to compute feasible box}
        \State $\mathbf{L} \gets [-\infty]^d, \quad \mathbf{U} \gets [+\infty]^d$
        \State $c \gets v$
        \While{$\text{parent}(c) \neq \text{NIL}$}
            \State $p \gets \text{parent}(c)$
            \State $f \gets g(p), \quad \theta \gets \theta(p)$
            \If{$c == \text{left}(p)$}
                \State $U_f \gets \min(U_f, \theta)$
            \Else
                \State $L_f \gets \max(L_f, \theta)$
            \EndIf
            \State $c \gets p$
        \EndWhile
        
        \LineComment{Step 2: Construct baseline routed vector}
        \State $\mathbf{z} \gets \mathbf{0}^d$
        \ForAll{constrained features $f \in \{1, \dots, d\}$}
            \If{$L_f > -\infty \land U_f < +\infty$}
                \State $z_f \gets (L_f + U_f) / 2$
            \ElsIf{$U_f < +\infty$}
                \State $z_f \gets U_f - 1.0$
            \ElsIf{$L_f > -\infty$}
                \State $z_f \gets L_f + 1.0$
            \EndIf
        \EndFor
        \State $\mathbf{x}_{\text{base}} \gets (\mathbf{z} \oslash \mathbf{s}) + \mathbf{o}$
        
        \LineComment{Step 3: Instantiate critical boundary triplets}
        \State $g \gets g(v), \quad \theta_v \gets \theta(v)$
        \State $x^* \gets (\theta_v / s_g) + o_g$
        \ForAll{$x' \in \{ x^*, \text{prev}_{32}(x^*), \text{next}_{32}(x^*) \}$}
            \State $\mathbf{x}_{\text{probe}} \gets \mathbf{x}_{\text{base}}$
            \State $\mathbf{x}_{\text{probe}}[g] \gets x'$
            \State $\mathcal{P} \gets \mathcal{P} \cup \{(\mathbf{x}_{\text{probe}}, v)\}$
        \EndFor
    \EndFor
\EndFor
\State \Return $\mathcal{P}$
\end{algorithmic}
\end{algorithm}

\subsection{Algorithm 2: Cache-Blocked Lockstep Traversal with Exact Early Exit}

While ONNX Runtime provides exceptional multi-threaded execution, embedded systems operating in low-power or restart modes require an agile, pure-NumPy fallback engine. Traditional sequential tree traversals implemented via Python-level loops suffer from severe interpreter overhead ($O(M \cdot \text{depth})$ bytecode dispatches per row).

Algorithm~\ref{alg:lockstep_early_exit} introduces a two-tier traversal optimization:
\begin{enumerate}
    \item \textbf{Cache-Blocked Lockstep Traversal:} Instead of advancing one tree across all samples, or one sample down all trees, trees are grouped into chunks $C_k$ sized to fit into L1/L2 data cache ($\text{PAIRS\_PER\_BLOCK} \le 8192$ sample-tree pairs). All active rows within a block advance simultaneously down the tree levels using vectorized NumPy masking operations, collapsing interpreter loop count from $O(M \cdot \text{depth})$ to $O(\text{depth} \cdot \lceil M / |C_k| \rceil)$.
    \item \textbf{Exact Margin-Bounded Early Exit:} For binary classification where leaf values represent log-odds margins, let $m_t^{\min}$ and $m_t^{\max}$ denote the global minimum and maximum margins obtainable in tree $t$. At step $k$, let the accumulated margin for sample $i$ be $m_i$, and let $R^{\min} = \sum_{t \notin \bigcup_{j=1}^k C_j} m_t^{\min}$ and $R^{\max} = \sum_{t \notin \bigcup_{j=1}^k C_j} m_t^{\max}$ denote the remaining extremal margin bounds.
\end{enumerate}

\begin{algorithm}[t]
\caption{Cache-Blocked Lockstep Early Exit}
\label{alg:lockstep_early_exit}
\footnotesize
\begin{algorithmic}[1]
\Require Batch $\mathbf{X} \in \mathbb{R}^{B \times d}$, schedule $\mathcal{S}$, bounds $[m_t^{\min}, m_t^{\max}]$
\Ensure Predictions $\mathbf{y} \in \{0, 1\}^B$, tree counts $\mathbf{u} \in \mathbb{N}^B$
\State $\mathbf{Z} \gets (\mathbf{X} - \mathbf{o}) \odot \mathbf{s}$
\State $\mathbf{m} \gets \mathbf{0}^B, \quad \mathbf{u} \gets \mathbf{0}^B$
\State $\mathcal{A} \gets \{1, 2, \dots, B\}$ \Comment{Active flow indices}
\State $R^{\min} \gets \sum_{t=1}^M m_t^{\min}, \quad R^{\max} \gets \sum_{t=1}^M m_t^{\max}$

\ForAll{\text{chunk } $C_k \in \mathcal{S}$}
    \If{$\mathcal{A} == \emptyset$}
        \State \textbf{break}
    \EndIf
    \LineComment{Cache-blocked parallel leaf traversal}
    \State $\mathbf{m}[\mathcal{A}] \gets \mathbf{m}[\mathcal{A}] + \sum_{t \in C_k} \text{LockstepLeaves}(\mathbf{Z}[\mathcal{A}], t)$
    \State $\mathbf{u}[\mathcal{A}] \gets \mathbf{u}[\mathcal{A}] + |C_k|$
    \State $R^{\min} \gets R^{\min} - \sum_{t \in C_k} m_t^{\min}$
    \State $R^{\max} \gets R^{\max} - \sum_{t \in C_k} m_t^{\max}$
    
    \LineComment{Exact Early Exit Check}
    \ForAll{$i \in \mathcal{A}$}
        \If{$m_i + R^{\min} > 0$}
            \State $y_i \gets 1$ \Comment{Provably attack}
            \State $\mathcal{A} \gets \mathcal{A} \setminus \{i\}$
        \ElsIf{$m_i + R^{\max} < 0$}
            \State $y_i \gets 0$ \Comment{Provably benign}
            \State $\mathcal{A} \gets \mathcal{A} \setminus \{i\}$
        \EndIf
    \EndFor
\EndFor
\ForAll{$i \in \mathcal{A}$}
    \State $y_i \gets \mathbb{I}(m_i \ge 0)$ \Comment{Residual evaluations}
\EndFor
\State \Return $\mathbf{y}, \mathbf{u}$
\end{algorithmic}
\end{algorithm}

\begin{theorem}[Exact Parity Guarantee of Early Exit]
If $m_i + R^{\min} > 0$, sample $i$ will be classified as an attack ($y_i = 1$) under full evaluation regardless of subsequent tree paths. Conversely, if $m_i + R^{\max} < 0$, sample $i$ is guaranteed benign ($y_i = 0$). Terminating sample $i$ upon satisfying either condition incurs zero label divergence ($0\,\text{flip}$) from the complete ensemble.
\end{theorem}
\begin{proof}
The final score of sample $i$ across all $M$ trees is $S_i = m_i + \sum_{t \text{ unvisited}} \Delta_t$. Since $m_t^{\min} \le \Delta_t \le m_t^{\max}$ by definition, it follows that $m_i + R^{\min} \le S_i \le m_i + R^{\max}$. If $m_i + R^{\min} > 0$, then $S_i > 0$ strictly holds, yielding an attack verdict under the zero-threshold decision rule. Symmetrically, if $m_i + R^{\max} < 0$, then $S_i < 0$ strictly holds, yielding a benign verdict. Hence, early termination cannot alter the final sign of $S_i$.
\end{proof}
